% Part: proof-theory
% Chapter: normalization
% Section: segments

\documentclass[../../../include/open-logic-section]{subfiles}

\begin{document}

\olfileid{pt}{nor}{seg}

\olsection{في المقاطع والقطوع}

إذا تعقبت قاعدةُ~!!a{elimination} قاعدةَ~!!^a{introduction} كان ذلك
التفافًا في~!!a{proof}، وتطبيع~!!a{proof} حذف أمثال هذه الالتفافات. لكن
قاعدتي~\Elim{\lor} و\Elim{\lexists} تزيدان تعريف «الالتفاف المراد حذفه»
تعقيدًا؛ لأن قاعدة~!!{elimination} لا يلزم أن تعقب قاعدة~!!{introduction}
\emph{مباشرةً}، بل ربما توسطتهما~\Elim{\lor} أو~\Elim{\lexists}، مثلًا:
\[
\AxiomC{}
\DeduceC{$\lexists[x][!C(x)]$}
\AxiomC{$\Discharge{!A}{x}$}
\DeduceC{$!B$}
\DischargeRule{\Intro{\lif}}{x}
\UnaryInfC{$!A \lif !B$}
\DischargeRule{\Elim{\lexists}}{y}
\BinaryInfC{$!A \lif !B$}
\AxiomC{}
\DeduceC{$!A$}
\RightLabel{\Elim{\lif}}
\BinaryInfC{$!B$}
\DisplayProof
\]
وقد يقع بين~\Intro{\lif} و\Elim{\lif} أي عدد من تطبيقات~\Elim{\lor}
أو~\Elim{\lexists}، مثلًا:
\[
\AxiomC{}
\DeduceC{$\lexists[x][!C(x)]$}
\AxiomC{}
\DeduceC{$!D \lor !E$}
\AxiomC{$\Discharge{!A}{x}$}
\DeduceC{$!B$}
\DischargeRule{\Intro{\lif}}{x}
\UnaryInfC{$!A \lif !B$}
\AxiomC{$\Discharge{!A}{x}$}
\DeduceC{$!B$}
\DischargeRule{\Intro{\lif}}{x}
\UnaryInfC{$!A \lif !B$}
\RightLabel{\Elim{\lor}}
\TrinaryInfC{$!A \lif !B$}
\DischargeRule{\Elim{\lexists}}{y}
\BinaryInfC{$!A \lif !B$}
\AxiomC{}
\DeduceC{$!A$}
\RightLabel{\Elim{\lif}}
\BinaryInfC{$!B$}
\DisplayProof
\]
ويظهر من المثال أن تطبيق~\Elim{\lif} الواحد قد يكون في التفافات عدة، إذ
يقع بعد أكثر من تطبيق واحد لـ~\Intro{\lif}. ولكي نستوعب هذه الوجوه نعرف
الالتفافات بأنواع مخصوصة من متتاليات وقائع~!!{formula} في برهان
\Log{N1i}، ونسميها \emph{مقاطع}. والمقطع في مثالنا متتالية من وقائع
$!A \lif !B$، تتبع فيها نتيجة~$\Elim{\lor}$ أو~$\Elim{\lexists}$
مقدمتها الصغرى. وفيه مقطعان من هذا الضرب، في كل منهما ثلاث وقائع من
$!A \lif !B$. والحد الدقيق هو الآتي.

\begin{defn}
\emph{المقطع} في~!!{proof}~$\delta$ من~\Log{N1i} متتالية الوقائع
$!A_1$، \dots، $!A_n$ لـ~!!{formula}~$!A$ بعينها في~$\delta$، وتخضع
للشروط الآتية:
\begin{enumerate}
  \item لا تكون~$!A_1$ نتيجة~\Elim{\lor} أو~\Elim{\lexists}؛
  \item لا تكون~$!A_n$ مقدمةً صغرى لـ~\Elim{\lor} أو~\Elim{\lexists}؛
  \item لكل~$i = 1$، \dots، $n-1$، تكون~$!A_i$ مقدمةً صغرى لـ
  \Elim{\lor} أو~\Elim{\lexists}، وتكون~$!A_{i+1}$ نتيجة ذلك الاستدلال.
\end{enumerate}
وعدد~$n$ هو \emph{طول} هذا المقطع.
\end{defn}

وليس كل مقطع التفافًا يراد حذفه؛ والمقاطع التي يراد حذفها هي
\emph{القطوع}.

\begin{defn}
\emph{القطع} مقطع تكون~$!A_n$ فيه المقدمة الكبرى لاستدلال~!!a{elimination}؛
وإما أن يكون~$n=1$، وتكون~$!A_1=!A_n$ نتيجة استدلال~!!a{introduction}
أو~\FalseInt، وإما أن يكون~$n>1$.
\end{defn}

ولا يشترط أن يبدأ مقطع القطع بنتيجة قاعدة~!!a{introduction}؛ وإنما
يشترط أن ينتهي بالمقدمة الكبرى لقاعدة~!!a{elimination}. لكن المقطع الذي
طوله~$1$ لا يعد قطعًا إلا إذا كان نتيجة قاعدة~!!{introduction} أو
\FalseInt.

\begin{defn}
\emph{رتبة} القطع هي~$\depth{!A}$. وأما \emph{رتبة القطع}
$\cutrank{\delta}$ لـ~!!a{proof}~$\delta$ فهي أعظم رتب قطوع~$\delta$،
ونجعلها~$0$ إذا خلا من القطوع. ونسمي القطع \emph{أقصى} في~$\delta$ إذا
كانت رتبته~$\cutrank{\delta}$، أي بلغت الغاية الممكنة. وأما
\emph{طول القطع} لـ~$\delta$ فمجموع أطوال قطوعه القصوى كلها، ونجعله~$0$
إذا لم يوجد قطع.
\end{defn}

\end{document}
